% =============================================================================
%  The Unitary Singularity:
%  A Formal Rejection of the Multiverse via
%  Algorithmic Parsimony
%
%  Author: Rafael D. De Paz
%  Date:   March 2026
%  Target: Foundations of Physics / Studies in HPS
% =============================================================================
\documentclass[12pt, a4paper]{article}

% --- Packages ---
\usepackage[utf8]{inputenc}
\usepackage[T1]{fontenc}
\usepackage{amsmath, amssymb, amsthm, mathtools}
\usepackage{geometry}
\usepackage{hyperref}
\usepackage{cleveref}
\usepackage{enumitem}
\usepackage{graphicx}
\usepackage{booktabs}
\usepackage{microtype}
\usepackage{natbib}

% --- Page Geometry ---
\geometry{margin=1in}

% --- Theorem Environments ---
\theoremstyle{plain}
\newtheorem{theorem}{Theorem}[section]
\newtheorem{lemma}[theorem]{Lemma}
\newtheorem{proposition}[theorem]{Proposition}
\newtheorem{corollary}[theorem]{Corollary}

\theoremstyle{definition}
\newtheorem{definition}[theorem]{Definition}
\newtheorem{assumption}[theorem]{Assumption}
\newtheorem{axiom}[theorem]{Axiom}

\theoremstyle{remark}
\newtheorem{remark}[theorem]{Remark}

% --- Custom Commands ---
\newcommand{\R}{\mathbb{R}}
\newcommand{\N}{\mathbb{N}}
\newcommand{\K}{\mathcal{K}}

% --- Hyperref Setup ---
\hypersetup{
  colorlinks=true,
  linkcolor=blue,
  citecolor=blue,
  urlcolor=blue,
  pdftitle={The Unitary Singularity},
  pdfauthor={Rafael D. De Paz}
}

% =============================================================================
\begin{document}

% --- Title ---
\title{
  \textbf{The Unitary Singularity: \\
  A Formal Rejection of the Multiverse \\
  via Algorithmic Parsimony}
}
\author{
  Rafael D. De Paz \\
  \textit{Independent Researcher} \\
  \texttt{me@rdepaz.com}
}
\date{March 2026}
\maketitle

% =============================================================================
% ABSTRACT
% =============================================================================
\begin{abstract}
The Many-Worlds Interpretation (MWI) of quantum mechanics and the Eternal
Inflation model of cosmology both posit the existence of infinite,
non-interacting parallel universes. We subject these hypotheses to three
independent formal analyses: (1) Kolmogorov complexity, demonstrating that
the algorithmic description length of a single universe with non-deterministic
observers is exponentially shorter than that of infinite branching histories;
(2) the Principle of Least Action, showing that universal branching at every
quantum event violates the variational principle that governs all known
physics; and (3) information-theoretic testability, proving that causally
disconnected universes are formally indistinguishable from noise under any
physically realizable measurement protocol. We conclude that the multiverse
hypothesis fails the Principle of Computational Parsimony---it is not
\emph{wrong} in the sense of being contradicted by evidence, but rather
\emph{unfounded} in the sense that it is algorithmically equivalent to
invoking an infinite number of unobservable entities to explain a finite
set of observations. A single deterministic universe with non-deterministic
observers (the ``Unitary Singularity'') is the maximally parsimonious
ontology consistent with all known physics.
\end{abstract}

\vspace{1em}
\noindent\textbf{Keywords:} multiverse, Many-Worlds Interpretation,
Kolmogorov complexity, algorithmic parsimony, fine-tuning,
Occam's razor, quantum foundations.

\vspace{0.5em}
\noindent\textbf{2020 MSC:} 81P05, 68Q30, 00A30, 83F05.

% =============================================================================
% 1. INTRODUCTION
% =============================================================================
\section{Introduction}\label{sec:intro}

The multiverse hypothesis---the claim that our universe is one of infinitely
many---has become a prominent feature of contemporary physics and cosmology,
motivated primarily by:

\begin{enumerate}[label=(\alph*)]
  \item The \textbf{Fine-Tuning Problem}: The fundamental constants of
    physics appear tuned to permit complex structure and life
    \citep{Barrow1986, Rees2000}.
  \item The \textbf{Measurement Problem} in quantum mechanics: The MWI
    \citep{Everett1957} avoids wavefunction collapse by postulating that
    all branches of the state vector are equally real.
  \item \textbf{Eternal Inflation}: Inflationary cosmology generically
    produces pocket universes with varying physical constants
    \citep{Guth1981, Linde1986}.
\end{enumerate}

While these motivations are substantive, the multiverse hypothesis introduces
an extraordinary ontological commitment: the existence of infinitely many
unobservable entities. We argue that this commitment is not justified by the
explanatory work it performs.

\subsection{Main Results}

\begin{enumerate}
  \item \textbf{Kolmogorov Parsimony} (\S\ref{sec:kolmogorov}): The
    algorithmic description of a branching multiverse has Kolmogorov
    complexity $\K(M) \to \infty$, whereas the single-universe model
    has finite $\K(U) < \infty$.

  \item \textbf{Variational Violation} (\S\ref{sec:least_action}): Universal
    branching at every quantum event requires infinite energy and storage,
    violating the Principle of Least Action.

  \item \textbf{Information-Theoretic Untestability}
    (\S\ref{sec:testability}): Causally disconnected universes carry zero
    mutual information with our observational data; they are formally noise.

  \item \textbf{The Unitary Alternative} (\S\ref{sec:unitary}): A single
    deterministic universe with a non-deterministic observer is sufficient
    to explain all known phenomena, including fine-tuning and quantum
    indeterminacy.
\end{enumerate}

% =============================================================================
% 2. KOLMOGOROV PARSIMONY
% =============================================================================
\section{The Kolmogorov Argument}\label{sec:kolmogorov}

\begin{definition}[Kolmogorov Complexity]\label{def:kolmogorov}
  The Kolmogorov complexity $\K(x)$ of a string $x$ is the length of the
  shortest program that produces $x$ on a universal Turing machine
  \citep{Kolmogorov1965, LiVitanyi2008}:
  \begin{equation}
    \K(x) = \min_{p} \{ |p| : U(p) = x \}.
  \end{equation}
\end{definition}

\begin{theorem}[Parsimony Inequality]\label{thm:parsimony}
  Let $U$ denote the description of a single deterministic universe with
  $n$ observer-induced non-deterministic measurements, and let $M$ denote
  the description of a multiverse with branching at each measurement.
  Then:
  \begin{equation}\label{eq:parsimony}
    \K(U) = \K(\text{laws}) + O(\log n)
    \quad \ll \quad
    \K(M) = \K(\text{laws}) + O(2^n),
  \end{equation}
  since $M$ must specify the state of $2^n$ branches while $U$ specifies
  only the observed branch.
\end{theorem}

\begin{proof}
  The single-universe model requires: (i) the physical laws (finite
  description), and (ii) the outcomes of $n$ quantum measurements (at most
  $n$ bits, representable in $O(\log n)$ additional description via
  compression of random sequences). The multiverse model requires: (i) the
  same physical laws, plus (ii) the complete state specification of all
  $2^n$ branches. Since each branch is a full universe-state, $\K(M) \geq
  \K(\text{laws}) + 2^n \cdot \K(\text{state per branch})$. The inequality
  follows. \qedhere
\end{proof}

\begin{corollary}[Occam's Razor, Formalized]\label{cor:occam}
  By Solomonoff's theory of inductive inference \citep{Solomonoff1964},
  the prior probability of a hypothesis $h$ is $P(h) \sim 2^{-\K(h)}$.
  Therefore: $P(U) / P(M) \sim 2^{\K(M) - \K(U)} \to \infty$.
\end{corollary}

% =============================================================================
% 3. THE VARIATIONAL ARGUMENT
% =============================================================================
\section{The Variational Argument}\label{sec:least_action}

\begin{axiom}[Principle of Least Action]\label{ax:least_action}
  All known fundamental physics obeys the Principle of Least Action: the
  path taken by a system between two states is the one that extremizes
  the action $S = \int \mathcal{L}\, dt$ \citep{Feynman1942, Landau1976}.
\end{axiom}

\begin{proposition}[Branching Violates Parsimony of Action]\label{prop:branching}
  The MWI requires that at every quantum measurement, the \emph{entire
  universe} duplicates into $k$ branches (where $k$ is the dimension of
  the measurement basis). The total action of the multiverse after $n$
  measurements is:
  \begin{equation}
    S_{\text{total}} = \sum_{i=1}^{k^n} S_i,
  \end{equation}
  where $S_i$ is the action of the $i$-th branch. This sum grows
  exponentially, whereas the single-universe action $S_U$ is bounded.
  A system that minimizes action cannot also maximize the number of
  copies of itself.
\end{proposition}

\begin{remark}
  Proponents of MWI argue that branches do not require ``additional
  energy'' because the total wavefunction is unitary. However, unitarity
  governs the \emph{norm} of the state vector, not its \emph{ontological
  commitment}. The question is not whether the mathematics is consistent
  but whether the ontology is parsimonious.
\end{remark}

% =============================================================================
% 4. INFORMATION-THEORETIC TESTABILITY
% =============================================================================
\section{Information-Theoretic Testability}\label{sec:testability}

\begin{definition}[Mutual Information]\label{def:mi}
  The mutual information between two random variables $X$ (our universe's
  data) and $Y$ (a parallel universe's data) is:
  \begin{equation}
    I(X; Y) = H(X) + H(Y) - H(X, Y),
  \end{equation}
  where $H$ denotes Shannon entropy \citep{Shannon1948, Cover2006}.
\end{definition}

\begin{theorem}[The Noise Theorem]\label{thm:noise}
  If universe $Y$ is causally disconnected from universe $X$ (i.e., no
  physical signal can travel between them), then $X$ and $Y$ are
  statistically independent, and:
  \begin{equation}
    I(X; Y) = 0.
  \end{equation}
  A hypothesis that predicts entities with zero mutual information with
  all observational data is empirically indistinguishable from a
  hypothesis that predicts no such entities.
\end{theorem}

\begin{proof}
  Causal disconnection implies $P(X, Y) = P(X) P(Y)$. Then
  $H(X, Y) = H(X) + H(Y)$, and $I(X; Y) = 0$. By the likelihood
  principle, any observational data $D$ satisfies
  $P(D \mid \text{multiverse}) = P(D \mid \text{single universe})$,
  rendering the multiverse unfalsifiable. \qedhere
\end{proof}

% =============================================================================
% 5. THE UNITARY ALTERNATIVE
% =============================================================================
\section{The Unitary Singularity}\label{sec:unitary}

\begin{definition}[Unitary Singularity]\label{def:unitary}
  The \emph{Unitary Singularity} is the ontological position that:
  \begin{enumerate}[label=(\roman*)]
    \item There exists exactly one physical universe.
    \item Its global dynamics are deterministic (governed by unitary
      evolution of the state vector).
    \item Local non-determinism arises from the observer's epistemic
      limitations: incomplete knowledge of initial conditions, not
      ontological branching.
  \end{enumerate}
\end{definition}

\begin{theorem}[Sufficiency of the Unitary Singularity]\label{thm:sufficiency}
  The Unitary Singularity accounts for all empirically observed phenomena
  that motivate the multiverse:
  \begin{enumerate}[label=(\alph*)]
    \item \textbf{Fine-tuning:} Explained by the non-invertibility of the
      causal generator function $f: \mathcal{G} \to \mathcal{O}$, where
      $f^{-1}$ is undefined from within the observable frame. The
      constants are not ``selected'' from a distribution but are
      structural features of a unique system.
    \item \textbf{Quantum indeterminacy:} Born-rule probabilities arise
      from the observer's position within the system (self-locating
      uncertainty), not from ontological branching.
    \item \textbf{Wavefunction collapse:} The ``collapse'' is an update
      of the observer's state, not a physical process in the universe.
  \end{enumerate}
\end{theorem}

% =============================================================================
% 6. DISCUSSION
% =============================================================================
\section{Discussion}\label{sec:discussion}

\subsection{Comparison with Existing Critiques}

Our analysis complements but extends existing critiques of the multiverse:

\begin{center}
\begin{tabular}{lll}
  \toprule
  \textbf{Critique} & \textbf{Author(s)} & \textbf{Our Extension} \\
  \midrule
  Unfalsifiability & Popper (1959) & Formalized via $I(X;Y)=0$ \\
  Ontological excess & Occam/Lewis & Quantified via $\K(M) \gg \K(U)$ \\
  Action violation & --- (novel) & Proved via variational principle \\
  Observer sufficiency & Fuchs (QBism) & Unified with parsimony bounds \\
  \bottomrule
\end{tabular}
\end{center}

\subsection{Limitations}

\begin{enumerate}
  \item Our Kolmogorov analysis assumes a fixed universal Turing machine;
    the complexity values are machine-dependent up to an additive constant.
  \item The variational argument addresses ontological cost, not
    mathematical consistency---MWI remains internally consistent.
  \item The fine-tuning resolution via non-invertibility is an axiom,
    not a derivation.
\end{enumerate}

% =============================================================================
% 7. CONCLUSION
% =============================================================================
\section{Conclusion}\label{sec:conclusion}

We have subjected the multiverse hypothesis to three independent formal
tests:

\begin{enumerate}
  \item \textbf{Algorithmic parsimony:} $\K(U) \ll \K(M)$.
  \item \textbf{Variational principle:} Branching maximizes total action.
  \item \textbf{Information content:} $I(X; Y) = 0$ for disconnected
    universes.
\end{enumerate}

In each case, the multiverse fails: it is less parsimonious, violates the
variational principle, and carries zero information content beyond what a
single universe provides. The Unitary Singularity---one deterministic
universe with epistemic non-determinism---is the simplest ontology
consistent with all known physics.

The multiverse is not the ``other worlds.'' It is the garbage collection
of the wavefunction: variables that were considered but never executed.

% =============================================================================
% REFERENCES
% =============================================================================
\bibliographystyle{plainnat}
\bibliography{references}

\end{document}
